\capitulo{7}{Conclusiones y Líneas de trabajo futuras}

El resultado del TFG es un flujo funcional y desplegado, no solo un diseño de arquitectura. Una solicitud puede entrar desde el portal o desde la Logic App, pasa por la misma validación y clasificación, se conserva en Blob Storage y queda disponible para las consultas protegidas y las métricas. El script de verificación recorre estas operaciones sobre la URL pública de Azure.

Las decisiones más útiles surgieron al reducir el alcance. Flask encajó mejor que un framework completo porque la API tiene pocas rutas y no necesita un modelo relacional. Por la misma razón se eligió Blob Storage frente a SQLite o Azure SQL. El pipeline de GitHub Actions previsto al principio se descartó a favor de scripts PowerShell y Azure CLI que podían ejecutarse, revisar y demostrar desde el entorno disponible.

La parte de seguridad también dejó una conclusión concreta: trasladar una clave a Key Vault no basta si la aplicación necesita otra credencial para leerla. La identidad administrada resuelve esa relación entre App Service y Key Vault sin añadir secretos al repositorio. SonarCloud ayudó a revisar configuraciones de infraestructura y a distinguir entre vulnerabilidades reales, decisiones intencionadas y aspectos pendientes de una solución de producción.

El trabajo permitió practicar de forma conjunta los siguientes ámbitos:

\begin{itemize}
    \item Seguridad cloud.
    \item Gestión de secretos.
    \item APIs REST.
    \item Automatización de procesos.
    \item Arquitecturas escalables.
    \item Clasificación automática de información operativa.
\end{itemize}

El prototipo cumple los objetivos planteados, aunque mantiene límites claros. La autenticación Bearer protege consultas técnicas, pero no sustituye la identidad individual de usuarios; Blob Storage sirve para el volumen de prueba, pero no ofrece las consultas de una base de datos; y el clasificador es explicable y repetible, pero no aprende de casos anteriores. Estas limitaciones orientan las líneas de continuación.

Como líneas futuras se plantean:

\begin{itemize}
    \item Evaluar un servicio de lenguaje solo cuando exista un conjunto de solicitudes etiquetadas con el que comparar su precisión frente a las reglas actuales.
    \item Incorporar autenticación de usuarios con Azure AD.
    \item Crear dashboards avanzados de monitorización y alertas operativas.
    \item Ampliar el portal web con seguimiento de solicitudes para usuarios no técnicos.
    \item Migrar la persistencia a Azure SQL o Cosmos DB para consultas más complejas.
    \item Incorporar notificaciones automáticas por correo o Teams cuando cambie el estado de una solicitud.
\end{itemize}
